\section{Recovery, scale and invariant capacity}\label{sec:meaning}
\subsection{Recovery without cumulative native gain}\label{sec:account}
The returning sign in Theorem~\ref{thm:credit} has an exact storage meaning. Empty external preparation and~\eqref{eq:balances} give, for all future exchange times,
\begin{equation}\label{eq:cumulative}
H_N(0)-H_N(t)=H_C(t)+H_R(t)
=\int_0^tP_{NC}(s)\dd s\ge0.
\end{equation}
Thus $H_N(t)\le H_N(0)$ even during a later credit interval.
The positive increment~\eqref{eq:gain} recovers earlier supplied trial storage; it does not create cumulative native gain above the initial stock.

At the comparison witness,
\begin{equation}\label{eq:comparisonbackflow}
\widehat z^{(0)}(T)=f_0,\qquad
\widehat y^{(0)}(T)=-\frac m\omega f_0.
\end{equation}
Their negative dot product identifies receiver backflow in that comparison. The finite threshold proves exact total native credit and initial exact receiver activity. It does not separately prove $P_{CR}(t_*)<0$ for every coupling satisfying~\eqref{eq:threshold}; an assertion about that different power requires its own finite bound.

\subsection{Fixed scaled windows and transferred amounts}\label{sec:scales}
For any fixed finite $S>0$, scaling and~\eqref{eq:externalbounds} imply
\begin{equation}\label{eq:stockscale}
\begin{aligned}
H_C(\tau/\kappa)+H_R(\tau/\kappa)
&=\frac{\eta_N^2}{2\kappa^2}\norm{V_\kappa(\tau)}_{CR}^2\\
&\le\frac{\eta_N^2mF_{\max}^2S^2}{2\kappa^2},
\qquad 0\le\tau\le S.
\end{aligned}
\end{equation}
Taking $S=T$ covers launch through the witness; $S=5T/4$ also covers the later interval. With all other parameters fixed, the external stock is
$O(\eta_N^2/\kappa^2)$ on such scaled windows. Equation~\eqref{eq:powererror} similarly gives
\begin{equation}\label{eq:powerscale}
P_{NC}(t_*)=-\frac{\eta_N^2D}{\kappa}
+O\!\left(\frac{\eta_N^2}{\kappa^2}\right).
\end{equation}
These are large-coupling statements on fixed scaled windows, not long-time decay statements. They show that strict existence of credit can coexist with a small absolute transferred stock. A small stock does not determine the fraction of it recovered: that question requires a specified numerator, denominator and time window. The present theorem establishes neither a substantial returned share nor maximum accessible extraction.

\subsection{The unchanged invariant ceiling}\label{sec:capacity}
The predecessor's Theorem 5.1 preserves the complex bilinear quantity
\begin{equation}\label{eq:invariant}
Z_2(\Omega)=\sum_{i=1}^3\Omega_i^2,\qquad
I=|\Omega|^2,\qquad
H_N(t)\ge\frac{\alpha_N}{2}|Z_2(\Omega_0)|.
\end{equation}
For the initial phase span $\phi$, combining its net-extraction bound with~\eqref{eq:cumulative} yields
\begin{equation}\label{eq:capacity}
\begin{aligned}
0\le\frac{H_N(0)-H_N(t)}{H_N(0)}
&\le\frac{I_0-|Z_2(\Omega_0)|}{I_0}
\le1-\cos\phi\\
&<1-\cos(\pi/12)\simeq0.0340741737109317.
\end{aligned}
\end{equation}
The denominator is positive initial core trial storage. The approximately
$3.4074\%$ maximum-span ceiling is a bound on net extraction, not an achieved fraction, gross repeated work, receiver share or diversion percentage. Its persistence needs invariance of $Z_2$, not continued membership in the initial coherent domain. A later native gain is entirely compatible with this ceiling.
